% source: https://github.com/pfradin/luadraw/discussions/159#discussioncomment-15175370 (pfradin, block 1)
\documentclass[border=5pt]{standalone}
\usepackage[svgnames]{xcolor}
\usepackage[3d]{luadraw}
\usepackage{fourier-otf}
\begin{document}

\begin{luadraw}{name=parallelep}
local g = graph3d:new{ window={-8,8,-8,8},bbox=false, viewdir={-30,70} }
Hiddenlinestyle = "dashed";Hiddenlines = true

-- data
local a, b, h = 3, 6, 5
local BAC = 60
local EAC = 60
local EAB = 60
local A = Origin

-- Calculation of B, C and D
local B = -a*vecI
local C1 = rotate3d(B,BAC, {A,vecK})
local D1 = C1+B-A
local Delta = pt3d.dot(B-A,D1-B)^2-a^2*(a^2-b^2)
local D = B + (-pt3d.dot(B-A,D1-B)+math.sqrt(Delta))/a^2*(D1-B)
local C = D-B+A

-- Calculation of E 
local u, v = pt3d.normalize(C-A), pt3d.normalize(B-A) 
-- we assume then E-A = x*u+y*v+z*vecK
local uv = pt3d.dot(u,v)
local alpha, beta, det = h*math.cos(EAC*deg), h*math.cos(EAB*deg), 1-uv^2
local x, y = (alpha-beta*uv)/det, (beta-alpha*uv)/det
local norm = pt3d.abs2(x*u+y*v)
local z = math.sqrt(h^2-norm)
local E = A+x*u+y*v+z*vecK

-- drawing
local P = parallelep(A,B-A,C-A,E-A)
g:Dpoly(P, {color="orange", mode=mShadedHidden})
g:Dlabel3d("$A$",A,{pos="S"}, "$B$",B,{}, "$C$",C,{}, "$D$",D,{},"$E$",E,{pos="N"}, "$F$",E+B-A,{}, "$H$",D+E-A,{},"$G$",C+E-A,{})

-- verification in the terminal
print( angle3d(C-A,E-A)*rad, angle3d(B-A,E-A)*rad, angle3d(C-A,B-A)*rad, pt3d.abs(B-A), pt3d.abs(D-A), pt3d.abs(E-A) )

g:Show()
\end{luadraw}
\end{document} 
